
\section{Unified Archimedean Softmax: Stability Theorem and Proof}
\label{sec:uas_theorem}

This section presents the second main analytical result: positive Harris recurrence of the CTMC under the Unified Archimedean Softmax (UAS) routing law. UAS extends raw softmax to heterogeneous servers by incorporating service rates into a prior-weighted variational objective and employing a weighted Lyapunov function.

\subsection{UAS Routing Law}
\label{sec:uas:routing}

\begin{definition}[UAS potential]
\label{def:uas_potential}
For temperature $\temp > 0$, define the UAS potential
\begin{equation}
\label{eq:uas_potential}
\uaspot{i}(\Qstate, \srvrate) = \frac{\Qlen{i}+1}{\srvrate_i} - \frac{1}{\temp}\log \srvrate_i.
\end{equation}
\end{definition}

The corresponding routing law is
\begin{equation}
\label{eq:uas_law}
\rprob{i}(\Qstate, \srvrate) = \frac{\srvrate_i \exp\!\bigl(-\temp (\Qlen{i}+1)/\srvrate_i\bigr)}{\sum_{j=1}^{\Nservers} \srvrate_j \exp\!\bigl(-\temp (\Qlen{j}+1)/\srvrate_j\bigr)},
\end{equation}
equivalently, $\rprob{i} = e^{-\temp \uaspot{i}} / \sum_j e^{-\temp \uaspot{j}}$.

\subsection{Variational Form}
\label{sec:uas:variational}

The UAS routing law admits a variational characterization based on the Kullback--Leibler divergence~\cite{cover2006elements}.

\begin{proposition}[UAS weighted variational form]
\label{prop:uas_variational}
Let $\totcap = \sum_{i=1}^{\Nservers} \srvrate_i$, $\uasprior{i} = \srvrate_i / \totcap$, and $\uasenergy{i}(\Qstate, \srvrate) = (\Qlen{i}+1)/\srvrate_i$. Define the prior-weighted entropy-regularized objective
\begin{equation}
\label{eq:uas_variational_obj}
\mathcal{G}_{\mathrm{UAS}}(p) = \sum_{i=1}^{\Nservers} p_i\, \uasenergy{i}(\Qstate, \srvrate) + \frac{1}{\temp}\KL{p}{r}
\end{equation}
over $\Simplex$. Its unique minimizer is the UAS routing law~\eqref{eq:uas_law}, and the minimum value is
\begin{equation}
\label{eq:uas_min_value}
-\frac{1}{\temp}\log \sum_{i=1}^{\Nservers} \uasprior{i}\, e^{-\temp\, \uasenergy{i}(\Qstate, \srvrate)}.
\end{equation}
\end{proposition}

\begin{proof}
Expand the KL divergence:
\begin{equation}
\label{eq:kl_expansion}
\KL{p}{r} = \sum_i p_i \log p_i - \sum_i p_i \log \uasprior{i}
= \sum_i p_i \log p_i - \sum_i p_i \log \srvrate_i + \log \totcap.
\end{equation}
Substituting into~\eqref{eq:uas_variational_obj} shows that $\mathcal{G}_{\mathrm{UAS}}(p)$ is algebraically identical to the standard UAS entropy-regularized objective $\sum_i p_i \uaspot{i} + \frac{1}{\temp} \sum_i p_i \log p_i$, up to a constant shift of $+\frac{1}{\temp}\log\totcap$. Strict convexity of $\mathcal{G}_{\mathrm{UAS}}$ in $p$ follows from the strict convexity of $\KL{p}{r}$ over the open simplex, guaranteeing a unique minimizer. The first-order optimality conditions on the simplex yield
\begin{equation}
\label{eq:uas_foc}
\log p_i = \mathrm{const} - \temp\, \uasenergy{i}(\Qstate, \srvrate) + \log \uasprior{i}
= \mathrm{const} - \temp\, \frac{\Qlen{i}+1}{\srvrate_i} + \log \srvrate_i,
\end{equation}
which after normalization gives the UAS routing law~\eqref{eq:uas_law}. Substituting the minimizer back into the objective yields the weighted $\logsumexp$ value~\eqref{eq:uas_min_value}.
\end{proof}

\subsection{Structural Interpretation}
\label{sec:uas:interpretation}

The UAS routing law integrates three effects into its routing score:
\begin{itemize}
\item \textbf{Queue pressure} via $\Qlen{i}$: longer queues receive less routing mass.
\item \textbf{Service heterogeneity} via $\srvrate_i$: the energy variable $\uasenergy{i} = (\Qlen{i}+1)/\srvrate_i$ normalizes queue length by service speed, and the prior $\uasprior{i} = \srvrate_i/\totcap$ favors faster servers.
\item \textbf{Entropy regularization} via the softmax form: routing probabilities remain smooth and interior.
\end{itemize}

At low load the prior term dominates and UAS favors faster servers; under congestion the queue-dependent energy term redistributes mass away from heavily loaded servers. The transition between these regimes is smooth by construction of the softmax form.

\subsubsection{Empty-State Behavior}
If all queues are empty ($\Qlen{i} = 0$ for all $i$), the routing law reduces to
\begin{equation}
\label{eq:uas_empty_state}
p_i(\mathbf{0}, \srvrate) \propto \srvrate_i\, e^{-\temp/\srvrate_i}.
\end{equation}
At large $\temp$ (low temperature, exploitative regime), $e^{-\temp/\srvrate_i} \to 0$ for slow servers while remaining close to $1$ for fast servers, so UAS concentrates routing mass on the fastest server.
At small $\temp$ (high temperature, exploratory regime), $e^{-\temp/\srvrate_i} \to 1$ for all $i$, so UAS approaches proportional routing ($p_i \propto \srvrate_i$).
This smooth interpolation between exploitative and proportional routing at the idle state distinguishes UAS from threshold-based policies.

\subsubsection{Pairwise Crossover Law}
For any two servers $f$ and $s$, the routing probabilities are equal ($p_f = p_s$) when
\begin{equation}
\label{eq:uas_crossover_1}
\log \srvrate_f - \temp\frac{\Qlen{f}+1}{\srvrate_f} = \log \srvrate_s - \temp\frac{\Qlen{s}+1}{\srvrate_s}.
\end{equation}
Rearranging yields the crossover boundary at which routing mass transfers between the two servers:
\begin{equation}
\label{eq:uas_crossover_2}
\frac{\Qlen{f}}{\srvrate_f} - \frac{\Qlen{s}}{\srvrate_s} = \frac{\log \srvrate_f - \log \srvrate_s}{\temp} + \left(\frac{1}{\srvrate_s} - \frac{1}{\srvrate_f}\right).
\end{equation}
The right-hand side is a constant (for fixed $\temp$ and service rates), so the crossover condition defines a hyperplane in the normalized queue-length space $\Qlen{i}/\srvrate_i$.
For JSQ, the analogous crossover is $\Qlen{f} = \Qlen{s}$; for JSSQ, it is $\Qlen{f}/\srvrate_f = \Qlen{s}/\srvrate_s$.
UAS interpolates between these extremes through the $\log$-ratio term, which vanishes as $\temp \to \infty$ (recovering JSSQ-like behavior) and dominates as $\temp \to 0$ (favoring the faster server regardless of queue imbalance).
This provides structural intuition for why Calibrated UAS with $\calbeta < 1$ may improve performance: tempering the normalization $\srvrate_i^{\calbeta}$ shifts the crossover boundary toward the faster server.

\subsection{Lyapunov Argument}
\label{sec:uas:lyapunov}

For the heterogeneous stability proof, we use the weighted Lyapunov function
\begin{equation}
\label{eq:uas_lyapunov}
\lyap_{\mathrm{UAS}}(\Qstate) = \frac{1}{2}\sum_{i=1}^{\Nservers} \frac{\Qlen{i}^2}{\srvrate_i}.
\end{equation}

\begin{lemma}[UAS drift identity]
\label{lem:uas_drift}
Under the UAS routing law~\eqref{eq:uas_law}, the generator action on $\lyap_{\mathrm{UAS}}$ is
\begin{equation}
\label{eq:uas_drift_identity}
\gen \lyap_{\mathrm{UAS}}(\Qstate)
= \arrate \sum_{i=1}^{\Nservers} \rprob{i}(\Qstate, \srvrate)\,\frac{\Qlen{i} + 1/2}{\srvrate_i}
- \sum_{i=1}^{\Nservers} \Qlen{i}
+ \frac{1}{2}\sum_{i=1}^{\Nservers} \Ind_{\Qlen{i}>0}.
\end{equation}
\end{lemma}

\begin{proof}
The generator applied to $\lyap_{\mathrm{UAS}}$ yields, for each arrival to server~$i$, the increment $\frac{1}{2\srvrate_i}\bigl[(\Qlen{i}+1)^2 - \Qlen{i}^2\bigr] = \frac{\Qlen{i}+1/2}{\srvrate_i}$ (after simplification), weighted by rate $\arrate \rprob{i}$. For each departure from server~$i$ (rate $\srvrate_i$ when $\Qlen{i}>0$), the increment is $\frac{1}{2\srvrate_i}\bigl[(\Qlen{i}-1)^2 - \Qlen{i}^2\bigr] = \frac{-\Qlen{i}+1/2}{\srvrate_i}$. Summing over all servers gives~\eqref{eq:uas_drift_identity}.
\end{proof}

\begin{lemma}[UAS arrival-term bound]
\label{lem:uas_arrival}
Let $\uasenergy{i}(\Qstate, \srvrate) = (\Qlen{i}+1)/\srvrate_i$ and $\uasprior{i} = \srvrate_i/\totcap$. Under the UAS routing law,
\begin{equation}
\label{eq:uas_arrival_bound}
\sum_{i=1}^{\Nservers} \rprob{i}(\Qstate, \srvrate)\,\frac{\Qlen{i}+1/2}{\srvrate_i}
\leq \frac{\abs{\Qstate}_1 + \Nservers}{\totcap}.
\end{equation}
\end{lemma}

\begin{proof}
By the weighted variational identity (\cref{prop:uas_variational}),
\begin{equation}
\label{eq:uas_var_id}
\sum_{i=1}^{\Nservers} \rprob{i}\, \uasenergy{i} + \frac{1}{\temp}\KL{p}{r}
= -\frac{1}{\temp}\log \sum_{i=1}^{\Nservers} \uasprior{i}\, e^{-\temp\, \uasenergy{i}}.
\end{equation}
Since $\KL{p}{r} \geq 0$,
\begin{equation}
\label{eq:uas_kl_drop}
\sum_{i=1}^{\Nservers} \rprob{i}\, \uasenergy{i}
\leq -\frac{1}{\temp}\log \sum_{i=1}^{\Nservers} \uasprior{i}\, e^{-\temp\, \uasenergy{i}}.
\end{equation}
Apply Jensen's inequality to the convex map $u \mapsto e^{-\temp u}$:
\begin{equation}
\label{eq:uas_jensen}
\sum_{i=1}^{\Nservers} \uasprior{i}\, e^{-\temp\, \uasenergy{i}}
\geq \exp\!\left(-\temp \sum_{i=1}^{\Nservers} \uasprior{i}\, \uasenergy{i}\right).
\end{equation}
Therefore
\begin{equation}
\label{eq:uas_jensen_result}
-\frac{1}{\temp}\log \sum_{i=1}^{\Nservers} \uasprior{i}\, e^{-\temp\, \uasenergy{i}}
\leq \sum_{i=1}^{\Nservers} \uasprior{i}\, \uasenergy{i}
= \frac{1}{\totcap}\sum_{i=1}^{\Nservers} \srvrate_i \frac{\Qlen{i}+1}{\srvrate_i}
= \frac{\abs{\Qstate}_1 + \Nservers}{\totcap}.
\end{equation}
Hence
\begin{equation}
\label{eq:uas_x_bound}
\sum_{i=1}^{\Nservers} \rprob{i}(\Qstate, \srvrate)\,\frac{\Qlen{i}+1}{\srvrate_i} \leq \frac{\abs{\Qstate}_1 + \Nservers}{\totcap}.
\end{equation}
Finally,
\begin{equation}
\label{eq:uas_half_bound}
\sum_{i=1}^{\Nservers} \rprob{i}\,\frac{\Qlen{i}+1/2}{\srvrate_i}
= \sum_{i=1}^{\Nservers} \rprob{i}\,\frac{\Qlen{i}+1}{\srvrate_i}
- \frac{1}{2}\sum_{i=1}^{\Nservers} \frac{\rprob{i}}{\srvrate_i}
\leq \sum_{i=1}^{\Nservers} \rprob{i}\,\frac{\Qlen{i}+1}{\srvrate_i},
\end{equation}
which gives~\eqref{eq:uas_arrival_bound}.
\end{proof}

\subsection{Main Theorem and Proof}
\label{sec:uas:theorem}

\begin{theorem}[UAS positive Harris recurrence]
\label{thm:uas}
Assume $\totcap = \sum_{i=1}^{\Nservers} \srvrate_i > \arrate$. Under the UAS routing law~\eqref{eq:uas_law}, the CTMC is non-explosive, irreducible, and positive Harris recurrent.
\end{theorem}

\begin{proof}
Non-explosion and irreducibility follow by the same argument as in \cref{thm:softmax}: the total jump rate is bounded by $\arrate + \totcap$, and all rates are strictly positive. The weighted Lyapunov function $\lyap_{\mathrm{UAS}}$ defined in~\eqref{eq:uas_lyapunov} is norm-like on $\Nats^\Nservers$.

Define
\begin{equation}
\label{eq:uas_epsilon}
\varepsilon_{\mathrm{UAS}} = \frac{\totcap - \arrate}{\totcap}
\end{equation}
and
\begin{equation}
\label{eq:uas_R}
R_{\mathrm{UAS}} = \frac{\arrate \Nservers}{\totcap} + \frac{\Nservers}{2}.
\end{equation}

The drift identity (\cref{lem:uas_drift}) gives
\[
\gen \lyap_{\mathrm{UAS}}(\Qstate)
= \arrate \sum_{i=1}^{\Nservers} \rprob{i}\,\frac{\Qlen{i}+1/2}{\srvrate_i}
\;-\; \underbrace{\sum_{i=1}^{\Nservers} \Qlen{i}}_{= \abs{\Qstate}_1}
\;+\; \underbrace{\frac{1}{2}\sum_{i=1}^{\Nservers} \Ind_{\Qlen{i}>0}}_{\leq\, \Nservers/2}.
\]
Applying the arrival-term bound (\cref{lem:uas_arrival}) to the first sum and bounding the indicator sum by~$\Nservers/2$,
\begin{align}
\gen \lyap_{\mathrm{UAS}}(\Qstate)
&\leq \arrate \cdot \frac{\abs{\Qstate}_1 + \Nservers}{\totcap}
  \;-\; \abs{\Qstate}_1
  \;+\; \frac{\Nservers}{2} \nonumber \\
&= \frac{\arrate}{\totcap}\,\abs{\Qstate}_1
  + \frac{\arrate \Nservers}{\totcap}
  - \abs{\Qstate}_1
  + \frac{\Nservers}{2} \nonumber \\
&= -\!\underbrace{\left(1 - \frac{\arrate}{\totcap}\right)}_{=\,\varepsilon_{\mathrm{UAS}}}\, \abs{\Qstate}_1
  + \underbrace{\frac{\arrate \Nservers}{\totcap} + \frac{\Nservers}{2}}_{=\, R_{\mathrm{UAS}}}.
\label{eq:uas_final_drift}
\end{align}

For the compact set
\begin{equation}
\label{eq:uas_compact}
C_{\mathrm{UAS}} = \left\{\Qstate \in \Nats^\Nservers : \abs{\Qstate}_1 \leq \frac{R_{\mathrm{UAS}} + 1}{\varepsilon_{\mathrm{UAS}}}\right\},
\end{equation}
we have $\gen \lyap_{\mathrm{UAS}}(\Qstate) \leq -1$ for all $\Qstate \notin C_{\mathrm{UAS}}$. By the continuous-time Foster--Lyapunov criterion, the CTMC is positive Harris recurrent.
\end{proof}

\subsection{Properties of the UAS Bound}
\label{sec:uas:discussion}

Four structural properties of the UAS stability result merit emphasis.

\paragraph{Independence from raw softmax}
\label{sec:uas:discussion:independence}
UAS is \emph{not} a corollary of \cref{thm:softmax}. It uses a different variational objective (prior-weighted KL divergence $\KL{p}{r}$ instead of uniform entropy), a different Lyapunov function ($\lyap_{\mathrm{UAS}} = \frac{1}{2}\sum_i \Qlen{i}^2/\srvrate_i$ instead of $\frac{1}{2}\sum_i \Qlen{i}^2$), and a different arrival-term bound (Jensen's inequality on the prior-weighted $\logsumexp$ instead of the entropy ceiling $\entropy \leq \log \Nservers$). These are structurally distinct proof paths.

\paragraph{Role of the prior}
The service-rate prior $\uasprior{i} = \srvrate_i / \totcap$ is the key structural difference from raw softmax. It enables the KL-based variational identity~\eqref{eq:uas_var_id} and the Jensen bound~\eqref{eq:uas_jensen}, which together eliminate the $(\log \Nservers)/\temp$ additive penalty that appears in the raw softmax bound~\eqref{eq:softmax_arrival_bound}.

\paragraph{Temperature-independent stability guarantee}
The drift constants $\varepsilon_{\mathrm{UAS}} = (\totcap - \arrate)/\totcap$ and $R_{\mathrm{UAS}} = \arrate\Nservers/\totcap + \Nservers/2$ depend only on $(\arrate, \srvrate_1, \dots, \srvrate_\Nservers)$ and are \emph{independent of the temperature~$\temp$}.
This is a structural advantage over raw softmax, whose constant $R = (\arrate \log \Nservers)/\temp + C_1$ degrades as $\temp \to 0$; UAS is provably stable at the same rate for all $\temp > 0$.

\paragraph{Motivation for empirical extensions}
Although UAS is provably stable, its parameterization $(\Qlen{i}+1)/\srvrate_i$ is not necessarily optimal for empirical performance, motivating the broader Calibrated UAS family and the stability-certified calibrated subclass developed in \cref{sec:calibrated_uas}.

